\documentclass[10pt,a4paper]{article}
\usepackage[margin=0.8in]{geometry}
\usepackage{amsmath,amssymb}
\usepackage[hidelinks]{hyperref}
\setlength{\parindent}{0pt}
\setlength{\parskip}{5pt}
\newcommand{\head}[1]{\par\medskip\textbf{#1}\par}
\begin{document}
\begin{center}
{\large\bfseries Joined VSL stack: analytical review addendum}\\[3pt]
2026-10-21 --- JM1; no new traffic law or numerical experiment
\end{center}

\textbf{Meaning of the result.} The complete remaining-time objective is
more than a formal infinite rollout with a conditional chain rule. For this
fixed model it admits a common finite representation over the allowed
advice interval and has rational formulas (ratios of polynomials in advice)
between finitely many possible switching values. This does not make it globally smooth or establish
a useful information or optimizer advantage.

\head{1. Same model and assumptions}
Use the unchanged
\href{vsl_joined_candidate_contract_and_qualification.pdf}{joined contract},
sections 1--7: fixed $N=10$ rows, continuous stocks $0\leq n_{i\ell}\leq J=15$,
common FD, selected H budgets, normalized requests, positive equal nominal
priorities, capped sharing, one forward fallback, and one-row forward or
diagonal advancement. Final E can only merge into A's outlet. All movements
use old stock. The initial cohort is independent of advice; there are no
later arrivals or blocked downstream outlets. Advice is one constant
$v\in I=[30,60]$ km/h, with the stated time/space bounds. Empty donors skip
undefined ratios. The fixed member's preferences and budgets use finitely
many rational formulas, comparisons and min/max branches.

Let $h_i=N-i+1$, $M(n)=\sum_{i,\ell}n_{i\ell}$ and
\begin{equation}
 V(n)=\sum_{i,\ell}h_i n_{i\ell},\qquad
 \eta=\Delta x\rho_c=15/4,\qquad c(v)=v\Delta t/\Delta x\in[1/2,1].
 \label{eq:definitions}
\end{equation}
Here $V$ is mass-weighted remaining row count, not a speed. Every realized
movement, including an exit, reduces $V$ by its mass.

\head{2. The progress argument supplies a uniform finite prefix}
Regard desired fractions $b$ as independent elements of their fixed compact
allowed simplices, keeping terminal E-to-A compulsory. For an empty donor,
any such $b$ represents the same zero requests. Define $\Gamma(n,b,v)$ as
the sum of all realized boundary movement amounts in one update. The
budget/allocation laws give
\begin{equation}
 V(n^+)=V(n)-\Gamma(n,b,v),\qquad M(n)\leq V(n)\leq NM(n).
 \label{eq:progress}
\end{equation}
The expanded map $\Gamma$ is continuous in $(n,b,v)$: budgets and requests
are continuous, capped sharing has unique continuous flow amounts at demand
caps and zero supply, and fallback takes minima of continuous residuals.
This statement does not assume continuity of the native incentive map $b(n)$.

If $V>0$, take the furthest occupied row. Its downstream row is empty or
is the free outlet. Each allowed destination there has positive receiving
budget, each positive donor has positive sending, and its requests sum to
that sending. Positive priorities therefore serve some positive mass.
Hence $\Gamma>0$ whenever $V>0$, for every admissible $b$.

On the compact set of physical stocks, allowed $b$, $v\in I$ and $V\geq\eta$,
continuity and positivity give a minimum $\mu>0$. As $V$ has a finite uniform
upper bound, \eqref{eq:progress} implies a finite common bound on steps to
$V<\eta$, independently of the native branch sequence or advice. Then
$M<\eta$; no inflow can undo this. The existing contract's free-tail proof
applies: all requests are served, each donor advances $c(v)n_{i\ell}$, and
the remaining value is exact. Choose one integer $K_*$ beyond this common
entry bound. For every $v\in I$,
\begin{equation}
 \mathcal J(v)=\Delta t\sum_{k=0}^{K_*-1}M(n^k(v))
              +\frac{\Delta x}{v}V(n^{K_*}(v)).
 \label{eq:finite}
\end{equation}
No numerical estimate of $\mu$ or simulated trajectory is used here. This
is an existence argument, not an efficient clearance algorithm or a claim
of exact finite emptying.

\newpage
\head{3. Finite branch formulas imply regular advice intervals}
Fix the admitted initial state. On a chosen branch, density lookahead is
linear, incentive ratios are rational wherever evaluated, geometry terms
are constant, H budgets are min/max choices of rational terms, and donor
normalization is rational. For capped sharing, fix fully served requests
$H$ and entitlement-limited requests $U$. In the constrained branch,
\begin{equation}
 \lambda=\frac{r-\sum_{a\in H}m_a}{|U|},\qquad
 z_a=m_a\ (a\in H),\quad z_a=\lambda\ (a\in U),\quad |U|>0.
 \label{eq:sharing}
\end{equation}
The unconstrained/zero cases have their explicit finite formulas. Fallback
and stock updates preserve rationality. There are finitely many choices per
update and $K_*$ is finite, so each candidate trace gives rational states,
branch conditions and full value \eqref{eq:finite} as functions of $v$.

Every nonidentically-zero rational comparison has finitely many roots and
denominator zeros on $I$: its numerator and denominator are polynomials of
finite degree. Identically zero comparisons are persistent ties, not an
infinite collection of separate switches; undefined-ratio branches are
excluded by their own conditions. Taking the finite union over candidate
traces yields a finite exceptional set. On each remaining open interval,
the selected formulas are fixed and $\mathcal J$ is rational with nonzero
denominators, hence smooth. Values at exceptional points may jump or have
unequal one-sided slopes. No global minimum-attainment claim follows.

For such an interval, write the selected step formula as $F^{[k]}$ and
$s^k=\partial n^k/\partial v$. The contract's derivative is justified:
\begin{align}
 s^0&=0,\qquad
 s^{k+1}=F_n^{[k]}s^k+F_v^{[k]},\nonumber\\
 \mathcal J_v&=\Delta t\sum_{k<K_*}\mathbf1^\top s^k
       +\frac{\Delta x}{v}h^\top s^{K_*}
       -\frac{\Delta x}{v^2}V(n^{K_*}). \label{eq:derivative}
\end{align}
The three value terms concern prefix traffic, residual-state sensitivity
and direct residual travel time. A persistent state tie need not admit a
global state Jacobian: the selected smooth formula along the advice curve
is what is differentiated. Ordinary derivatives can still be zero. A finite
partition need not be cheaply enumerable. No AD/FD verification, smoothing,
continuous-time saltation or continuum-limit result is asserted.

\head{4. A bounded negative information result in the free regime}
If $M(n^0)<\eta$, the entire episode already lies in the exact-tail regime
for every allowed $v$. Therefore
\begin{equation}
 \mathcal J(v,n^0)=\frac{\Delta x}{v}V(n^0),\qquad
 \mathcal J_v=-\frac{\Delta x}{v^2}V(n^0).
 \label{eq:free}
\end{equation}
For any nonempty cohort, maximum advice uniquely minimizes this expression.
If two initializations share row totals but differ in lane allocation, they
share $V$ and the whole objective for every advice. Hidden lane allocation
then has zero action value. The empty cohort has zero value for all advice.

The existing X/Z contrast has $M_0=38.25>\eta$ and therefore is not covered
by this sufficient free-regime condition. Its different local services do
not prove different full-objective optimum advice. Outside the free regime,
fixed physical receiving/ceiling budgets and state-dependent preferences
prevent treating advice as a universal time rescaling, but no general
positive information-gap or maximum-advice dominance theorem is proved here.

\textbf{Review boundary.} These are deductions from the selected finite,
closed-cohort member. Continuing inflow, blocked exits, changed action timing,
heterogeneous-FD mapping or different traffic laws need their own qualification.
No source calibration, numerical gradient correctness, optimizer advantage,
physical transfer, full-lane necessity or discrete-order necessity follows.
The earlier model and its equations remain unchanged.
\end{document}
